% Artículo VII: pruebas preexistentes incorporadas al cuerpo.
% Propietario íntegro: fuentes_conservadas/17_rebote_einstein_cartan.tex.
% Incluye la ley de escala y las pruebas de rebote y persistencia.
% La amplitud s_0 es la salida de la realización material antecedente;
% este fragmento no certifica por sí solo la evaluación de esa realización.

\section{Dilución de memoria y persistencia del rebote}
\label{vii:sec:dilucion-rebote}

La evolución de la memoria y la expansión espacial son dos lecturas de
una misma sucesión de estados. El índice de refinamiento puede eliminarse
entre ambas; esa eliminación determina el exponente de dilución que la
realización cosmológica transmite a su contribución torsional.
La amplitud inicial depende de la evaluación de la corriente material
concreta en su respuesta variacional. La ley de escala y el teorema
de rebote que siguen conservan esa dependencia.

El teorema~\ref{vii:thm:amplitud-material-explicita} determina esa
amplitud como funcional del estado material y del volumen comóvil.
El corolario~\ref{vii:cor:dominio-positivo-conservado} proporciona
un dominio explícito de estados con amplitud positiva y conservada.
En ese dominio, la evaluación de \(s_0\) precede a la integración
cosmológica que se desarrolla a continuación. La densidad ordinaria
\(\rho_0>0\) y su ley barotrópica se especifican adicionalmente en la
sección~\ref{vii:sec:umbral-rebote}: el testigo de positividad de la
amplitud no evalúa por sí mismo esa densidad material.

\subsection{Ley de dilución \(a^{-6}\)}
\label{vii:sec:dilucion}

\begin{proposition}[Eliminación del índice de memoria]
\label{vii:prop:dilucion}
Sean \(a_{\mathrm{ref}}>0\), \(\varphi>1\) y los lectores de la ruta
\[
 \frac{a_n}{a_{\mathrm{ref}}}=\varphi^{n/3},\qquad
 M_n=\varphi^{-2n}.
\]
Entonces
\[
 M_n=\left(\frac{a_n}{a_{\mathrm{ref}}}\right)^{-6}
\]
para todo índice de la sucesión. Si la realización de la densidad
cuadrática es \(s_n^2=s_*^2M_n\), con \(s_*^2>0\), se obtiene
\[
 s_n^2=s_*^2\left(\frac{a_n}{a_{\mathrm{ref}}}\right)^{-6}.
\]
\end{proposition}

\begin{proof}
Elevando la primera igualdad a la potencia \(-6\) se obtiene
\(\varphi^{-2n}\), que es exactamente el segundo lector. Multiplicar
por la amplitud inicial reproduce la segunda afirmación. El exponente
queda determinado por el cociente entre los exponentes de ambos
lectores; la amplitud es la evaluación del estado de referencia.
\end{proof}

En la realización homogénea continua se conserva esta ley constitutiva:
la densidad torsional se escribe \(s_0a^{-6}\), donde la normalización
del factor \(a\), el coeficiente material y la amplitud de referencia
se han absorbido en \(s_0\). La ley de escala y la evaluación de
\(s_0\) son operaciones sucesivas de la construcción.

\subsection{Umbral homogéneo y transversalidad}
\label{vii:sec:umbral-rebote}

Consideremos la coordenatización homogénea, espacialmente plana, en unidades \(c=1\),
con \(G>0\) constante, densidades positivas \(\rho_0,s_0\) y parámetro
barotrópico constante \(w<1\). La realización especificada satisface
\(\Lambda_{\mathrm{eff}}=0\) y las ecuaciones
\begin{align}
 \rho_{\mathrm m}(a)&=\rho_0a^{-3(1+w)},&
 \rho_{\mathrm s}(a)&=s_0a^{-6},\\
 H^2&=\frac{8\pi G}{3}
        \bigl(\rho_{\mathrm m}-\rho_{\mathrm s}\bigr)
        =:F_0(a),&
 \dot H&=-4\pi G
        \bigl((1+w)\rho_{\mathrm m}-2\rho_{\mathrm s}\bigr).
 \label{vii:eq:dinamica-homogenea}
\end{align}
Los acoplamientos pertenecen a la sección dimensional ya construida.
La expresión de \(\rho_{\mathrm s}\) recoge el signo de su contribución
en esta realización gravitatoria.

\begin{theorem}[Existencia y unicidad del umbral transversal]
\label{vii:thm:rebote}
Existe un único \(a_{\mathrm b}>0\) con \(F_0(a_{\mathrm b})=0\):
\[
 a_{\mathrm b}^{\,3(1-w)}=\frac{s_0}{\rho_0}.
\]
La región permitida por \(H^2\geq0\) es \(a\geq a_{\mathrm b}\).
En el umbral, con
\(\rho_{\mathrm b}=\rho_{\mathrm m}(a_{\mathrm b})
                  =\rho_{\mathrm s}(a_{\mathrm b})>0\),
\[
 \dot H_{\mathrm b}=4\pi G(1-w)\rho_{\mathrm b}>0.
\]
La solución que atraviesa \(H=0\) presenta un mínimo estricto de \(a\),
con expansión local
\[
 a(t)=a_{\mathrm b}\left[
 1+2\pi G(1-w)\rho_{\mathrm b}(t-t_{\mathrm b})^2
 +O((t-t_{\mathrm b})^3)\right].
\]
\end{theorem}

\begin{proof}
La factorización
\[
 F_0(a)=\frac{8\pi G}{3}a^{-6}
             \bigl(\rho_0a^{3(1-w)}-s_0\bigr)
\]
reduce sus ceros al de una función estrictamente creciente de
\((0,\infty)\) sobre \((-s_0,\infty)\). Esto demuestra existencia,
unicidad y el signo de \(F_0\) a ambos lados del umbral.
La segunda ecuación de \eqref{vii:eq:dinamica-homogenea}, evaluada
en la igualdad de densidades, da la fórmula de \(\dot H_{\mathrm b}\).
Como \(\dot a=aH\), en el umbral
\(\ddot a_{\mathrm b}=a_{\mathrm b}\dot H_{\mathrm b}>0\).
Las ecuaciones son suaves para \(a>0\), por lo que su desarrollo de
Taylor da la expansión indicada. El cambio de signo de \(H\) es
transversal: pasa de contracción a expansión.
\end{proof}

Para \(w=1\), ambos términos tienen el mismo exponente y la igualdad de
amplitudes determina un caso degenerado. Para \(w>1\), el signo de
\(\dot H\) en una coincidencia de densidades es negativo.
El dominio \(w<1\) del teorema expresa exactamente el orden entre los
dos exponentes de dilución.

\subsection{Persistencia local bajo perturbaciones controladas}
\label{vii:sec:persistencia}

La persistencia del umbral utiliza dos estimaciones independientes:
el control del signo en los extremos y el de la derivada en el
intervalo. Sea \(I=[a_-,a_+]\subset(0,\infty)\), con
\(a_-<a_{\mathrm b}<a_+\), elegido de forma que
\[
 F_0(a_-)<0<F_0(a_+),\qquad
 m_I:=\inf_{a\in I}F_0'(a)>0.
\]
Tal intervalo existe porque
\[
 F_0'(a_{\mathrm b})
   =\frac{8\pi G}{a_{\mathrm b}}(1-w)\rho_{\mathrm b}>0.
\]
Definimos
\(\delta_I=\min\{-F_0(a_-),F_0(a_+)\}>0\).

\begingroup\fontsize{10}{11.7}\selectfont
\setlength{\parskip}{2pt}\setlength{\abovedisplayskip}{5pt}\setlength{\belowdisplayskip}{5pt}
\begin{theorem}[Persistencia local del rebote transversal]
\label{vii:thm:persistencia}
Sea \(R\in C^1(I)\) una perturbación que satisface
\[
 \|R\|_{\infty,I}<\delta_I,\qquad
 \|R'\|_{\infty,I}<m_I.
\]
Entonces \(F_R=F_0+R\) tiene un único cero \(a_R\in(a_-,a_+)\) y
\(F_R'(a_R)>0\). La realización \(H^2=F_R(a)\) admite el rebote local
transversal, con
\[
 \dot H_R=\frac{a_R}{2}F_R'(a_R)>0
\]
en el umbral. Si \(R(a,\eta)\) es conjuntamente \(C^1\), los ceros
persistentes dependen localmente de manera \(C^1\) del parámetro \(\eta\).
\end{theorem}

\begin{proof}
La primera cota conserva
\(F_R(a_-)<0<F_R(a_+)\); el teorema del valor intermedio proporciona
un cero interior. La segunda da
\(F_R'(a)\geq m_I-\|R'\|_{\infty,I}>0\) en \(I\), por lo que el cero
es único y simple.

En \(a>a_R\) cercano al cero,
\(F_R(a)=F_R'(a_R)(a-a_R)+o(a-a_R)\).
La integral
\[
 t-t_{\mathrm b}=\int_{a_R}^{a}
          \frac{d\xi}{\xi\sqrt{F_R(\xi)}}
\]
es finita y estrictamente creciente. Su inversa define la rama
expansiva; el reflejo temporal define la rama contractiva.
Ambas se unen con \(\dot a(t_{\mathrm b})=0\) y
\[
 a(t)-a_R
 =\frac{a_R^2F_R'(a_R)}4(t-t_{\mathrm b})^2
   +o((t-t_{\mathrm b})^2).
\]
Fuera del umbral, la derivación de \(H^2=F_R(a)\) da
\(\dot H=aF_R'(a)/2\); su límite continuo da la expresión positiva
en \(a_R\). La regularidad \(C^1\) de \(F_R\) garantiza la continuidad
de la aceleración en esta unión. Finalmente,
\(\partial_aF_R(a_R)>0\) permite aplicar el teorema de la función
implícita cuando la perturbación depende conjuntamente de \((a,\eta)\).
\end{proof}

La primera cota controla la existencia; la segunda asegura unicidad
local y transversalidad. La estabilidad demostrada corresponde a
ese umbral en el intervalo seleccionado. La amplitud torsional
interviene en su localización, mientras que la relación entre
exponentes determina la orientación temporal del cruce.

\par\endgroup
